\chapter{CONCLUSION}
\label{ch:conclusion}
\section{Self-Similar Voice as Task Assistance}
As spoken agents guide users through virtual environments, the voice that delivers their assistance becomes a design choice in its own right. Follow My Voice asks whether a synthetic voice resembling the participant changes search efficiency or experience compared with a neutral voice the participant has selected. Because resemblance, liking, and usefulness need not coincide, the study examines them separately, in a task where spoken guidance progressively reveals the intended object.

The proposed within-participant comparison uses two counterbalanced 32-trial blocks after one practice trial in the first assigned voice. Active completion time is the primary outcome. Retrospective workload, motivation, and agent-experience ratings, together with a debrief after both blocks, provide complementary evidence. Perceived-similarity ratings assess whether the synthesized voice achieves the intended resemblance, while completion-time measures evaluate search performance.

Gaze-contingent guidance supplies the shared context in which participants encounter the two voices. Local spatial reasoning, decisions from the local MiniLM model, explicit right-trigger selection, and traceable speech delivery support a controlled and inspectable comparison. These components are intended to make the voice comparison interpretable. This thesis does not estimate their independent effects or compare decision providers.

This thesis contributes a rationale and study design for examining self-similar voice, together with an implemented apparatus and measurement plan. Confirmatory participant results are not yet available. Before the study can estimate a voice-condition contrast, sample-size planning, pilot acceptance criteria, and alignment of the consent documents with cloud processing must be completed.

\section{Status of the Research Questions}
\textbf{RQ1} asks whether voice self-similarity impacts performance and perceived workload. It has not yet been answered empirically: no confirmatory contrast in active completion time is estimated here, and NASA-TLX responses have not been collected. \textbf{RQ2} asks whether voice self-similarity impacts eye-tracking indicators of workload and search efficiency. The measurement pipeline runs end to end on headset recordings with both voices (Section~\ref{sec:development-eye-measures}), but the detectors still require validation and no voice comparison is made. \textbf{RQ3} asks whether voice self-similarity impacts engagement in the search task; the engagement measure is defined, but no participant ratings have been collected. The hypotheses that the self-similar voice improves performance and lowers workload (H1), makes search more efficient (H2), and makes the task more engaging (H3) therefore remain to be tested.

For guidance delivery, the platform and its data dictionary allow sensing, state, decisions, and software playback to be inspected separately. The October 20, 2026 verification run summarized in Appendix~\ref{app:development} documents specific software checks. It provides neither a held-out guidance error rate nor evidence that all delayed confirmations have been eliminated. These records describe the supporting apparatus and its current limits.

\section{Next Empirical Step}
Before confirmatory recruitment, the study configuration will be fixed, the consent documents reconciled with the services actually used, the full procedure validated with wearers, and the primary analysis and sample-size justification finalized. The resulting data will record counterbalanced order, actual voice exposure, observation coverage, and the block referenced by each retrospective survey set. Only then can the results framework be replaced with estimates and the abstract and conclusion revised around observed evidence.

The intended contribution is evidence for choosing a guiding voice: whether a self-similar voice changes search efficiency in either direction, alters the experience of assistance, introduces discomfort, or leaves important differences unresolved. Those answers will be limited to the tested task and voice configurations.
